\section{Ultraproducts, gradings, and finite cells}
\label{sec:foundations}

Fix the integer $n\geq 1$ for which we seek a finite realization.
This section provides a rational algebraic description of finite
constructions in an ultraproduct of ordinary $p$-local spectra.  The
description keeps track of sphere degrees that vary with $p$.  In
particular, it will allow us to interpret a finite sequence of exterior
algebra extensions as cofiber sequences of spectra.

All categories, tensor products, limits, and colimits are understood
homotopically.  The notation $\map$ denotes a mapping spectrum and
$\Map=\Omega^\infty\map$ its mapping space.  A commutative algebra is
an $E_\infty$-algebra unless a strict differential graded model is
specified.  Chain degrees are homological: the differential, denoted
$\partial$, lowers degree by one.  For a discrete abelian group
$\Gamma$, a $\Gamma$-graded object is a family indexed by $\Gamma$;
its tensor product is convolution,
\[
 (U\otimes V)_\eta
   =\bigoplus_{\alpha+\beta=\eta}U_\alpha\otimes V_\beta.
\]
Formal weights introduce no symmetry signs.  In a chain model the
symmetry sign is determined only by chain degree.

\subsection{The ordinary categorical ultraproduct}

Choose a nonprincipal ultrafilter $\U$ on the set of odd primes.
A set belonging to $\U$ will be called \emph{large}.  Define
\begin{equation}\label{eq:found-ultraproduct}
 \cC=\prod_{\U}\Sp_{(p)},\qquad
 D=\prod_{\U}\Z_{(p)},\qquad
 \varpi=[(p)_p]\in D.
\end{equation}
Here the category in Equation~\eqref{eq:found-ultraproduct} is
\[
 \cC=\colim_{U\in\U}\prod_{p\in U}\Sp_{(p)},
\]
with transition functors given by restriction to smaller large sets.
Thus a sequence represents an object up to restriction to a large set.
We use successively larger universes when taking this colimit and its
Ind-completion.  This is a colimit of categories; no presentability or
preservation of infinite primewise colimits is part of the definition.

\begin{proposition}\label{found:ultraproduct}
The category $\cC$ is stable and symmetric monoidal, with finite
stable constructions and tensor products computed primewise.  If
$X=[(X_p)]$ and $Y=[(Y_p)]$, then
\begin{equation}\label{eq:found-mapping-spectrum}
 \map_{\cC}(X,Y)
  \simeq\colim_{U\in\U}\prod_{p\in U}
                   \map_{\Sp_{(p)}}(X_p,Y_p).
\end{equation}
Every mapping spectrum is rational, and the endomorphism algebra of
the tensor unit is $HD$.  The ring $D$ contains $\Q$, and $\varpi$ is
a nonzerodivisor.

Objects, arrows, and homotopies forming a diagram indexed by a finite
simplicial set can be represented simultaneously on a large set of
primes.  The same holds for specified finite cofiber diagrams.  Sequence
objects become compact in $\Ind(\cC)$.
\end{proposition}

\begin{proof}
Products and filtered colimits of stable categories preserve the
stable structure.  To see the mapping formula and the assertion about
finite diagrams directly, use simplicial presentations of the
categories.  Inner horn fillers and any prescribed finite diagram
involve only finitely many simplices.  Such data therefore pass through
the filtered colimit.  Mapping spaces are obtained from finite limit
tests on these presentations, and filtered colimits of spaces commute
with finite limits.  This gives the mapping-space version of
Equation~\eqref{eq:found-mapping-spectrum}; applying it with every
fixed integer suspension gives the assertion for mapping spectra.
The same mapping-space tests show that primewise finite limits and
colimits have the required universal properties.  The tensor product
and its coherent finite-arity structure pass through the construction
as well.  These facts also follow from the general categorical
ultraproduct results in the author version of
\cite[Proposition~3.12, Corollary~3.16, and Lemma~3.17]{BSS2020}.

For every fixed nonzero integer $m$, multiplication by $m$ is an
equivalence at every prime not dividing $m$.  It is consequently an
equivalence on each mapping spectrum in
Equation~\eqref{eq:found-mapping-spectrum}.  These spectra are therefore
rational.  Homotopy groups commute with products of spectra and with
filtered colimits, so their groups are the corresponding group
ultraproducts.  Each fixed positive stable stem of the sphere is finite;
this is the stable consequence of Serre's finiteness theorem
\cite[Chapter~V, \S3, Proposition~3]{Serre1951}.  Its $p$-localization
therefore vanishes for all but finitely many primes.  Negative stable
stems vanish, and the zeroth stem gives $D$.  Hence the endomorphism
algebra of the unit is the discrete commutative ring spectrum $HD$.

The inverse of any fixed nonzero integer exists in $D$, giving the
inclusion $\Q\subset D$.  If $\varpi a=0$, choose a sequence
representing $a$.  On a large set one has $p a_p=0$ in the domain
$\Z_{(p)}$, so $a_p=0$ there.  Thus $a=0$.  Finally, compactness of
the sequence objects in the Ind-completion is part of the defining
Yoneda embedding into that completion.
\end{proof}

The same rationality argument applies to ultraproducts of categories
of spectra with naive actions of finite groups that may vary with
$p$: multiplication by a fixed integer prime to $p$ is still an
equivalence.  We often omit $H$ when writing a discrete scalar ring
as a ring spectrum.  Notice that rationality in Proposition~\ref{found:ultraproduct}
concerns fixed integer degrees.  It does not discard homotopy in the
prime-dependent degrees used below.

\subsection{Rational algebra models and specified nullhomotopies}

\begin{lemma}\label{found:rectification}
Let $R$ be a discrete commutative ring containing $\Q$, and let
$\Gamma$ be a discrete abelian group.  Commutative algebras in
$\Gamma$-graded $HR$-module spectra admit strict commutative differential
graded $R$-algebra models, with no boundedness requirement on chain
degree.  This description is compatible with diagrams and with
derived base change under specified scalar algebras.
\end{lemma}

\begin{proof}
The symmetric monoidal equivalence between $HR$-module spectra and
the derived category of $R$-modules is
\cite[Theorem~7.1.2.13]{LurieHA}.  The commutative differential graded
algebra model and its comparison with $E_\infty$-algebras require
$\Q\subset R$, precisely our hypothesis
\cite[Propositions~7.1.4.10--7.1.4.11]{LurieHA}.
For clarity, the comparison uses a resolution of the commutative
operad by an $R$-linear $E_\infty$-operad.  Since the order of each
finite symmetric group is invertible in $R$, its coinvariant functor
is exact.  Comparison of free algebras is therefore a quasi-isomorphism;
the filtration by free cells extends this comparison to cellular
resolutions.  With the additional $\Gamma$-grading the same argument
is applied weight by weight, summing the weights of the inputs.  Under
a prescribed scalar algebra one first chooses a cofibrant model of
that algebra and then uses derived base change.  All these comparisons
are equivalences of categories with their mapping spaces, so they also
apply to diagrams and their specified homotopies.
\end{proof}

We will need more than the equivalence class of an algebra that kills
a cycle: its universal nullhomotopy must induce the intended map of
underlying cofibers.  The following formulation records that choice.

\begin{lemma}\label{found:free-null}
Let $B$ be a $\Gamma$-graded commutative differential graded algebra
over a discrete ring $R\supset\Q$, and let $x\in B_\eta$ be a cycle
of chain degree $r$.  The free extension
\[
 B//x:=B[z],\qquad
 |z|_{\mathrm{chain}}=r+1,\quad \wt(z)=\eta,\quad
 \partial z=x
\]
models the derived algebra pushout universally adjoining a null of
$x$.  In contrast, write $B/x$ for the underlying module cofiber of
multiplication by $x$, with its degree and weight shifts.

When $r=0$, there is an identification of underlying $B$-modules
\begin{equation}\label{eq:found-null-cofiber}
 B//x\simeq B/x.
\end{equation}
It uses the specified universal null $z$.  Consequently, an algebra
map from $B//x$ defined by a map from $B$ and a specified null of the
image of $x$ induces the corresponding map of module cofibers using
that null multiplied by the base input.
\end{lemma}

\begin{proof}
Let $F$ be the free commutative $R$-algebra on a cycle $u$ of degree
$r$ and weight $\eta$, and map $u$ to $x$.  The required pushout is
$B\otimes_F R$, where $F\to R$ sends $u$ to zero.  Resolve this
augmentation by the free disk algebra
\[
 F[z],\qquad \partial z=u.
\]
The complex spanned by $u,z$ is contractible.  Every positive symmetric
power of this disk is acyclic, since tensor powers are contractible
and symmetric-group coinvariants are exact over $R$.  Thus the disk
algebra maps by a quasi-isomorphism to $R$.
As an $F$-module it has an increasing filtration by the number of
copies of $z$, whose successive quotients are free shifted
$F$-modules.  This filtration computes the derived tensor with $B$.
After base change the result is exactly $B[z]$ with $\partial z=x$.
The argument respects the specified weight throughout.

If $r=0$, the variable $z$ is odd, hence exterior.  The underlying
module is the two-term extension of $B$ by the weight-$\eta$, degree-one
copy of $B$, with differential given by multiplication by $x$.  This
is Equation~\eqref{eq:found-null-cofiber}.  To identify the chosen
null without an ambiguity, one can first perform this calculation
over the polynomial cycle algebra $F=R[u]$.  Its augmentation is the
cofiber of multiplication by $u$, and its discrete target has no
degree-one ambiguity in the null of the generator.  Tensoring to $B$
gives the displayed comparison through the universal null.  A map
out of the disk algebra sends $z$ to the chosen target null; on the
underlying two-term module its value on a base multiple of $z$ is
therefore that same null multiplied by the image of the base element.
\end{proof}

\begin{remark}\label{found:scalar-pushouts}
The free cycle and disk algebras over $R$ are obtained by extension
of scalars from those over $\Q$.  Thus the cell in
Lemma~\ref{found:free-null} also has its stated universal property
when considered merely as an attachment over rational spectra.
More generally, a pushout of a span of algebras under common scalars
is also the pushout of that span without separately imposing those
scalars: the map from the middle algebra already supplies the common
scalar map and its compatibility.  We will use this observation for
power operations that can change the action of $D$.  Comparisons
involving specified nulls are transported together with their
specified homotopies along equivalences.
\end{remark}

\subsection{Coherent sphere gradings}

We next assign varying sphere degrees to a fixed lattice of formal
weights.  The required multiplicative coherence is imposed in the
rational category $\cC$.

\begin{lemma}\label{found:picard-grading}
Let $\mathcal D$ be a stable symmetric monoidal category with rational
mapping spectra, and let $\Gamma$ be a finite free abelian group with
a chosen basis.  Suppose the desired invertible object for each basis
element is a tensor square of an invertible object.  These objects
extend to a symmetric monoidal grading of $\mathcal D$ by the
discrete group $\Gamma$.

For two such gradings, prescribed equivalences on the basis extend
to a monoidal equivalence of gradings.  The construction can be made
relative to any specified subset of the basis.
\end{lemma}

\begin{proof}
The Picard space of $\mathcal D$ is a grouplike $E_\infty$-space and
hence is the infinite loop space of a connective spectrum $U$;
the recognition statement is
\cite[Remark~5.2.6.26]{LurieHA}.  A monoidal grading by $\Gamma$ is
equivalently a spectrum map $H\Gamma\to U$.  For $k\geq2$,
\[
 \pi_k U\cong\pi_{k-1}\map_{\mathcal D}(\one,\one),
\]
so $\tau_{\geq2}U$ is rational.

Put $Q_s=\cofib(S\to H\Z)$.  The rationalization of $S\to H\Z$
is an equivalence.  Moreover $Q_s$ is 2-connective, and
$\pi_2 Q_s=\pi_1^s=\Z/2$
\cite[Corollaries~4J.3--4J.4]{Hatcher2002}. Since maps from a rationally zero spectrum
to a rational spectrum vanish, there is an equivalence
\[
 \map(Q_s,U)\simeq\map(Q_s,\tau_{\leq1}U).
\]
Connectivity and truncation imply that its homotopy groups in degrees
at least zero vanish.  In degree minus one its only possible
contribution is
\[
 \pi_{-1}\map(Q_s,\tau_{\leq1}U)
   \cong\Hom(\Z/2,\pi_1 U),
\]
which is killed by two.  The fiber sequence
\[
 \map(Q_s,U)\longrightarrow\map(H\Z,U)
                    \longrightarrow\map(S,U)=U
\]
therefore shows that evaluation is an equivalence of spaces onto a
union of components of $\Omega^\infty U$.  The obstruction to a
component lying in that union takes values in a group killed by two.
Every double in $\pi_0 U$ lies in the image, so a specified tensor
square extends to a grading on $\Z$.

The equivalence onto a union of components also identifies all path
spaces between eligible objects.  Hence specified equivalences of
the generator objects lift to equivalences of gradings, with their
higher compatibilities.  Finally, $H\Gamma$ is the finite direct sum
of copies of $H\Z$ indexed by the basis.  The space of its maps to
$U$ is the corresponding product.  Applying the preceding argument
factor by factor proves both the assertion and its relative form.
\end{proof}

Define the lattices and additive functions
\begin{equation}\label{eq:found-weights}
 \begin{gathered}
  \Lambda=\langle w_0,\ldots,w_{n-1}\rangle,\qquad
  \Omega=\langle w_0,\ldots,w_n\rangle,\\
  d=d_p:\Omega\longrightarrow\Z,\qquad
  d(w_h)=d_h=2(p^h-1),\qquad q=d_1=2(p-1),\\
  \nu:\Omega\longrightarrow\Z,\qquad \nu(w_h)=1.
 \end{gathered}
\end{equation}
The symbols $w_h$ and their relations are independent of $p$, whereas
$d_h$ and $q$ depend on $p$.  Each sequence of spheres
$(S^{d_p(w_h)})_p$ is a tensor square.  By
Lemma~\ref{found:picard-grading} we may choose a monoidal sphere
grading $S(\eta)$ on $\Omega$ with
\[
 S(\eta)\simeq[(S^{d_p(\eta)})_p].
\]
Since $d_p(w_0)=0$, choose this grading first on
$\Omega/\langle w_0\rangle$ and then pull it back.  The $w_0$
direction is thus the unit direction.  Fixed homological suspensions,
which may be odd, are kept separate from these even sphere degrees.

\subsection{Normalized mapping objects}

For $\Gamma=\Omega$ or $\Lambda$ and $X\in\cC$, define
\begin{equation}\label{eq:found-normalized-maps}
 h_\Gamma(X)_\eta=\map_{\cC}(S(\eta),X),
 \qquad \eta\in\Gamma.
\end{equation}
For a primewise $\Gamma$-graded system $(X_p(\eta))$, the same
notation means
\[
 h_\Gamma(X)_\eta
   =\map_{\cC}\bigl(S(\eta),[(X_p(\eta))_p]\bigr).
\]
An ungraded algebra can also be regarded as constant in the weights,
using its multiplication for the weight-additive pairings.

\begin{lemma}\label{found:mapping-pairings}
The constructions $h_\Gamma$ are exact componentwise and possess
coherent lax symmetric monoidal pairings.  On algebraic inputs they
give graded $D$-algebras and their modules.  The pairings are balanced
over $D$ whenever its action is supplied by trivial sphere scalars.
They also respect additional diagram indices, including filtrations.
\end{lemma}

\begin{proof}
Mapping out of a fixed sphere is exact.  Given two input maps, tensor
them and use the chosen equivalence
$S(\alpha)\otimes S(\beta)\simeq S(\alpha+\beta)$; then compose
with the target pairing.  This defines the pairing at the pair of
weights $(\alpha,\beta)$.  Monoidality of the sphere grading and
separate exactness of tensor make these maps coherent at every finite
tuple of weights.  Their sums give convolution pairings in the
category of graded mapping spectra.  This construction requires no
infinite coproduct in $\cC$ itself.  The endomorphism algebra of the
unit acts centrally through tensor, so the pairings are $D$-balanced.
If the targets carry another diagram index, the same construction
commutes with its structure maps and its additive pairings.
\end{proof}

\begin{remark}\label{found:whole-systems}
For graded inputs we use an ultraproduct of categories of \emph{whole}
primewise graded systems.  Evaluation at any fixed index gives a
functor into $\cC$.  Primewise graded pairings then give the diagrams
and pairings used in Lemma~\ref{found:mapping-pairings}, compatibly
with restriction to large sets.  One may first pull these systems
back along prime-dependent additive maps from a common indexing
lattice.  This only uses a functor from the ultraproduct of whole
systems to systems in $\cC$; it does not identify these categories
or assert that every diagram in $\cC$ has such a lift.  The same
distinction applies to group actions, filtrations, and module
structures.  The simultaneous coherent sphere grading was chosen
only in $\cC$.  An individual sphere or cell shift used primewise is
simply the sphere of the indicated degree.
\end{remark}

Put $A=h_\Omega(\one)$.  By Lemma~\ref{found:rectification} we may
work with a strict commutative differential graded $D$-algebra model
of $A$.  Choose that model on the smaller grading
$\Omega/\langle w_0\rangle$ before pulling it back.  The resulting
model is strictly $w_0$-periodic.  Let
\[
 u_0\in A_{w_0},\qquad |u_0|_{\mathrm{chain}}=0,
\]
be the homogeneous periodic unit corresponding to $1\in A_0$.
We now show that finite algebraic constructions over $A$ describe
actual finite constructions from these sphere objects.

\subsection{Finite-cell comparison and the objects \texorpdfstring{$Y_j$}{Yj}}

For a graded module write $V\langle\eta\rangle_\theta
=V_{\theta-\eta}$ for weight shift.  A \emph{finite-cell} graded
$A$-module is obtained from finitely many homological suspensions of
the modules $A\langle\eta\rangle$ by finite sums and cofibers.  We
do not include retracts in this terminology.  Let $\cC_{\mathrm{cell}}$
be the corresponding full stable subcategory of $\cC$ generated by
the spheres $S(\eta)$.

\begin{proposition}\label{found:finite-cell-morita}
The functor $h_\Omega$ induces a symmetric monoidal equivalence
between $\cC_{\mathrm{cell}}$ and the category of finite-cell graded
$A$-modules.  For $X,Y\in\cC$, the natural comparison
\begin{equation}\label{eq:found-tensor-comparison}
 h_\Omega(X)\otimes_A h_\Omega(Y)
       \longrightarrow h_\Omega(X\otimes Y)
\end{equation}
is an equivalence whenever either input is in $\cC_{\mathrm{cell}}$.
These comparisons preserve units, compositions, and symmetric
multi-input products.  In particular, a commutative $A$-algebra whose
underlying module is finite-cell determines a commutative algebra
in $\cC$.
\end{proposition}

\begin{proof}
Invertibility gives
\[
 h_\Omega(S(\eta))_\theta
   \simeq\map_{\cC}(S(\theta-\eta),\one)
   =A_{\theta-\eta}.
\]
Thus a sphere generator maps to $A\langle\eta\rangle$.  For any
$X\in\cC$, mapping from this free graded module evaluates
$h_\Omega(X)$ at $\eta$, and therefore agrees with
$\map_{\cC}(S(\eta),X)$.  Exactness extends this comparison from
the generators to every finite-cell source.  It follows that
$h_\Omega$ is fully faithful on $\cC_{\mathrm{cell}}$.  A finite-cell
graded module is realized inductively: lift its attaching maps by
full faithfulness and take the corresponding cofibers in $\cC$.
This gives precisely the stated essential image.

The pairings of Lemma~\ref{found:mapping-pairings} are balanced over
$A$.  Their coherent balanced comparison is obtained from the tensor
bar construction: use the pairings with an arbitrary finite number
of inserted $A$-factors and then take its realization in graded
mapping spectra.  If $X=S(\eta)$,
Equation~\eqref{eq:found-tensor-comparison} is the weight-shift
equivalence supplied by invertibility.  Both sides are exact in
$X$, so the comparison remains an equivalence for every finite-cell
$X$, with $Y$ arbitrary.  Symmetry gives the other case.
The same bar pairings for several inputs and iterated relative
tensors retain compatibility with composition and symmetry.  Hence
the equivalence on the full tensor-closed finite-cell subcategories
is symmetric monoidal.  Commutative algebra objects in those full
subcategories transport across it, with all their coherence.
\end{proof}

The proposition holds with $\Lambda$ in place of $\Omega$, and with
the finite-cell subcategory generated by the spheres indexed by
$\Lambda$, by exactly the same generator and cofiber tests.  Only
finite-cell sources are involved; no conservativity assertion on
arbitrary objects is being made.

\begin{proposition}\label{found:koszul-realization}
Fix $0\leq r\leq n$. Suppose cycles $b_h$ of chain degree zero and
weight $w_h$, for $0\leq h\leq r$, have been chosen successively
in the strict algebras
\begin{equation}\label{eq:found-exterior-tower}
 \begin{gathered}
 A[e_0,\ldots,e_{j-1}],\qquad
 |e_h|_{\mathrm{chain}}=1,\quad \wt(e_h)=w_h,\\
 \partial e_h=b_h\in A[e_0,\ldots,e_{h-1}],
 \qquad b_0=\varpi u_0.
 \end{gathered}
\end{equation}
For $0\leq j\leq r+1$, these algebras determine commutative
algebras $Y_j\in\cC$, with $Y_0=\one$. For $0\leq j\leq r$,
the cycle $b_j$ represents a map
$S(w_j)\to Y_j$, and there is a cofiber sequence
\begin{equation}\label{eq:found-cofiber-tower}
 S(w_j)\otimes Y_j\xrightarrow{\,b_j\cdot-\,}Y_j
                    \longrightarrow Y_{j+1},
\end{equation}
whose second map is the algebra map adjoining the null $e_j$.
For $M=[(\MU_{(p)})_p]$ and $0\leq j\leq r+1$, there are
compatible equivalences
\begin{equation}\label{eq:found-mu-base-change}
 h_\Omega(M\otimes Y_j)
   \simeq h_\Omega(M)\otimes_A A[e_0,\ldots,e_{j-1}].
\end{equation}
\end{proposition}

\begin{proof}
Each variable $e_h$ is exterior.  The underlying $A$-module has the
finite basis of ordered exterior monomials $e_S$ for
$S\subseteq\{0,\ldots,j-1\}$.  Order these subsets by the largest
entry at which they differ.  In the differential of $e_S$, replacing
an index $h$ by indices strictly less than $h$ decreases this order;
terms with a repeated index vanish.  The differential therefore
provides a finite-cell filtration on this basis.  Proposition~\ref{found:finite-cell-morita}
transports the algebra and its unit to $Y_j$.

A degree-zero cycle at weight $w_j$ represents a class in
$\pi_0\map_{\cC}(S(w_j),Y_j)$.  Its product with the identity of
$Y_j$ corresponds to multiplication by $b_j$ on the algebraic module.
Lemma~\ref{found:free-null} identifies the extension by $e_j$, through
its specified null, with the cofiber of this multiplication map.
This proves Equation~\eqref{eq:found-cofiber-tower}, including its
unit and cofiber-map identifications.  Finally, $Y_j$ is finite-cell,
so Equation~\eqref{eq:found-tensor-comparison} with the other input
$M$ gives Equation~\eqref{eq:found-mu-base-change}, compatibly with
all the algebra maps and specified nulls.
\end{proof}

In particular, $Y_1$ is the cofiber of $\varpi$ on the unit.  The
remaining task is to construct the cycles $b_1,\ldots,b_n$ with the
required cobordism images.  The algebra structures above are available
in $\cC$ throughout that construction.  At the eventual passage to a
prime we will retain only finite diagrams of classes, products,
units, and cofiber sequences, as permitted by
Proposition~\ref{found:ultraproduct}.
